\documentclass[11pt,letterpaper]{article}
\usepackage[T1]{fontenc}
\usepackage{lmodern,amsmath,amssymb,amsthm,microtype}
\usepackage[margin=1in]{geometry}
\usepackage[hidelinks]{hyperref}
\providecommand{\tightlist}{\setlength{\itemsep}{0pt}\setlength{\parskip}{0pt}}
\DeclareUnicodeCharacter{2264}{\ensuremath{\le}}
\DeclareUnicodeCharacter{2265}{\ensuremath{\ge}}
\DeclareUnicodeCharacter{2260}{\ensuremath{\ne}}
\DeclareUnicodeCharacter{2212}{\ensuremath{-}}

\DeclareUnicodeCharacter{2211}{\ensuremath{\sum}}

\usepackage{longtable,booktabs,array}
\title{Excluding eighth weights below 136 in the nine-weight search below 150}
\author{Aafi Khan Malek\\ \small Independent researcher, Ahmedabad, India\\
\small\href{mailto:aafimalek2032@gmail.com}{aafimalek2032@gmail.com}}
\date{28 November 2026}
\begin{document}
\maketitle
\textbf{Theorem.} Let nine sorted positive real weights have pairwise
unit-separated subset sums. If \(149\le x_9<150\), then \(x_8\ge 136\).
The proof uses the preceding
\href{https://github.com/Aafimalek/novel-discoveries/tree/2d9416dca07cce2d85ba973076b788761f392f53/nine-separated-subset-sums/verification/rank-cuts-below-150/penultimate-cap-134}{cap-134 theorem} \(x_8\ge 135\) and a complete exact
replay of the next cap. This does not determine the optimum of the
nine-weight problem.

\section{Cap-specific implications}\label{cap-specific-implications}

The cap \(t=135\) forces coordinate-9 starts 146 through 172, so the
unsigned full-sign slots are \(q=38,\dots,64\). The magnitudes
\(\alpha_q=u_{a-q+1}\in[149-2q,151-2q)\) are consecutive, with common
\(64\le a\le74\) and \(b=148-a\). The established plus-two endpoints
then give the capacity ranges:

\begin{longtable}[]{@{}lr@{}}
\toprule\noalign{}
Half-open magnitude range & Capacity \\
\midrule\noalign{}
\endhead
\bottomrule\noalign{}
\endlastfoot
{[}1,21) & \(a-64\) \\
each {[}149−2q,151−2q), q from 64 down to 38 & 1 \\
{[}75,225) & 37 \\
{[}278,321) & \(148-2a\) \\
{[}373,10000) & \(a-20\) \\
\end{longtable}

The two wide gaps have width at least
\(\alpha_{38}-\alpha_{64}+4\ge56\). Flipping a sign at a weight strictly
below 28, or swapping signs at adjacent weights with a gap strictly
below 28, preserves membership at both cuts. The small-flip orbits on
complementary sign pairs give \(2^s\mid a,b\) for the number s of small
weights, hence \(2^s\mid148\) and \(s\le2\). A nonsmall coordinate
exists because otherwise every absolute signed magnitude is below
\(8\cdot28<373\). Thus \(x_3\ge28\). Expanded zero-based positions
{[}0,a), {[}a,b), and {[}b,128) retain their exact capacities; an empty
middle group for a=b permits no middle label.

For \(y=(129,\dots,136)\) write \(r_i=y_i-x_i\), \(R=\sum r_i\).
Balanced magnitudes at y are at most 16 and unbalanced ones at least
250. If any required form is unbalanced, \(R>175\) because its magnitude
is below 75. Otherwise \(\alpha_{38}\ge73\) is balanced and \(R\ge57\).
Thus \(R\ge57\) in either case, and \(r_3\le131-28=103\). The prior
theorem supplies \(x_8\ge135\), hence \textbf{\(r_8\le1\)}. For an
ordered deficit tail of length N with sum T, subtract \(r_8\) from each
entry and use the extreme prefix-sign averages of the first N−1 entries.
The signed tail is bounded by the corresponding endpoint expressions
\(\mu T+(\sum\sigma-N\mu)r_8\) over the T interval and
\(0\le r_8\le\min(1,T_{\rm upper}/N)\), rounding outward. No optional
additional tail cuts are used in the actual baseline.

\section{Deterministic coverage and exact
results}\label{deterministic-coverage-and-exact-results}

The 8,256 initial unit cells leave 7,360 before splitting by a. Of
80,960 two-deficit count profiles, 35,600 fail pure orbit size, 3,547
fail ordinary interval tests, and 19,181 fail the interval-conditioned
orbit assignment. The canonical residual is \textbf{22,632} profiles,
which the manifest partitions into \textbf{181} disjoint exhaustive
shards of at most 128 profiles.

All \textbf{181} clear with exact refinement: 109 at depth 3 and 72 at
depth 4. Summing the guarded worker logs gives:

\begin{center}
\small
\begin{tabular}{rrrrrr}
\toprule
Depth & Cells & Orbit size & Intervals & Orbit + intervals & Surviving \\
\midrule
3 & 1,119,917 & 249,659 & 142,052 & 711,965 & 16,241 \\
4 & 549,098 & 306,509 & 89,564 & 153,025 & 0 \\
\bottomrule
\end{tabular}
\end{center}

A lightweight independent audit checks the complete disjoint partition,
each source pin, and the upstream and downstream digests of each
completed depth. A separate arithmetic audit reconciles all 253 recorded
depths with their exact rejection logs and residual counts. The hashes
preserve the proof artifacts; they do not replace replaying the
arithmetic.

\section{Replay}\label{replay}

\begin{verbatim}
sha256sum -c SHA256SUMS
python3 parameters135.py
python3 coverage_audit.py
python3 arithmetic_totals.py
./replay-all.sh
\end{verbatim}

The last script copies the pinned source and checkpoint to a fresh
temporary directory, reruns phase 2 to demand canonical byte equality,
and recomputes all 181 shards serially with no reuse of old progress.
Each proof invocation is bounded by 35 seconds and 143,000 KiB virtual
memory, and memory and disk reserves are checked between shards. It uses
only the Python standard library and exact integer or fraction
arithmetic. The geometric implications are handwritten; this is not a
Lean theorem, a new upper witness, or an exclusion of any candidate with
\(x_8\ge136\).

\end{document}
